\documentclass[conference]{IEEEtran}
%\IEEEoverridecommandlockouts
% The preceding line is only needed to identify funding in the first footnote. If that is unneeded, please comment it out.
%Template version as of 6/27/2024

\usepackage{cite}
\usepackage{amsmath, amssymb, amsfonts}
\usepackage{algorithmic}
\usepackage{graphicx}
\usepackage{textcomp}
\usepackage{xcolor}

\usepackage{multicol}
\usepackage{amsmath}
\usepackage{amsfonts}
\usepackage{booktabs}
\usepackage{hyperref}

\def\BibTeX{{\rm B\kern-.05em{\sc i\kern-.025em b}\kern-.08em
    T\kern-.1667em\lower.7ex\hbox{E}\kern-.125emX}}

\begin{document}

\title{RSE-Supported Scientific Software Standardization: A Case Study of Eelbrain-BIDS}
% {\footnotesize \textsuperscript{*}Note: Sub-titles are not captured for https://ieeexplore.ieee.org  and
% should not be used}
%\thanks{Identify applicable funding agency here. If none, delete this.}
%}

\author{\IEEEauthorblockN{1\textsuperscript{st} Spencer Smith}
\IEEEauthorblockA{\textit{Computing and Software Department} \\
\textit{McMaster University}\\
Hamilton, Ontario, Canada\\
smiths@mcmaster.ca}
\and
\IEEEauthorblockN{2\textsuperscript{nd} Christian Brodbeck}
\IEEEauthorblockA{\textit{Computing and Software Department} \\
\textit{McMaster University}\\
Hamilton, Ontario, Canada}
\and
\IEEEauthorblockN{2\textsuperscript{nd} Zhangwenchi Li}
\IEEEauthorblockA{\textit{Computing and Software Department} \\
\textit{McMaster University}\\
Hamilton, Ontario, Canada}
}

\maketitle

\begin{abstract}
    Scientific software often starts from a relatively small project by a domain expert end user programmer.  In this nascent stage some existing domain standards for libraries, data structures, input formats or output formats may be neglected. As this software matures, the team may wish to adopt existing standards to improve interoperability, reproducibility and reusability. A research software engineer (RSE) may be asked to support standardization.  This can be challenging because of code coupling and the RSE not initially have the needed domain knowledge.  We propose the following steps for the RSE interacting with the domain expert: 1) Decide whether backward compatibility is desired, 2) Decide whether to design for potentially adopting a different standard in the future, 3) Learn necessary domain knowledge, 4) Understand and document the existing system architecture, 5) Design for the new standard, 6) Implement the change, 7) Verify the implementation, 8) Plan for migration for existing users of the software.  We have identified guidelines to support these steps, including incremental iterative development, information hiding, question-driven learning, multi-layered testing, and deprecation warnings.  This talk presents the process and guidelines via the example of the Eelbrain-BIDS project.  The Eelbrain-BIDS project aims to migrate the input dataset format of Eelbrain, a Python library for MEG (Magnetoencephalography) and EEG (Electroencephalography) data analysis, from a custom format to the neuroscience standard BIDS (Brain Imaging Data Structure).
\end{abstract}

\begin{IEEEkeywords}
research software, assurance case, verification and validation, medical image
analysis
\end{IEEEkeywords}

\section{Introduction}


\section{Background} \label{Sec_Background}


\section{Conclusions} \label{Sec_Conclusions}


\section*{Acknowledgements}

% \bibliographystyle{IEEEtran}
% \bibliography{../Resources/RSE-SSSS} 

\end{document}